\documentclass[10pt,UTF8]{article}

\usepackage{geometry}

\usepackage{ctex}

\geometry{a4paper,scale=0.9}
\ctexset{today=old}

\usepackage[bookmarksnumbered=true]{hyperref}  % 生成大纲


\usepackage{bm}
\usepackage{times}  % 设置正文字体为 Adobe Times Roman
% \usepackage{mathptmx} % 只加载数学字体
% \renewcommand{\rmdefault}{cmr} % 恢复正文字体为 Computer Modern Roman

% \usepackage{makecell}
% \usepackage{multirow} % 表格多行合并

% \usepackage{graphicx} %插入图片的宏包
% \usepackage{subfigure}
% \usepackage{float} %设置图片浮动位置的宏包

\usepackage{amsmath}
\usepackage{amssymb}
\usepackage{mathtools}

\usepackage{centernot}

\usepackage{physics}

% \usepackage{siunitx}
% % 关闭警告
% \ExplSyntaxOn
% \msg_redirect_name:nnn { siunitx } { physics-pkg } { none }
% \ExplSyntaxOff

%%%%%%%%%%%%%%%%%%%%%%%% tcolorbox %%%%%%%%%%%%%%%%%%%%%%%%%%
\usepackage[most]{tcolorbox}

\newenvironment{tipbox}
{\begin{tcolorbox}[colback=green!5!white,colframe=green!75!black]}
{\end{tcolorbox}}

\newenvironment{conceptbox}
{\begin{tcolorbox}[colback=blue!5!white,colframe=blue!75!black]}
{\end{tcolorbox}}

\newenvironment{thmbox}
{\begin{tcolorbox}[colback=yellow!5!white,colframe=yellow!75!black]}
{\end{tcolorbox}}

\newenvironment{definition box}
{\begin{tcolorbox}[colback=blue!5!white,colframe=blue!75!black,parbox=false,before upper=\par,before lower=\par]}
{\end{tcolorbox}}

\newenvironment{theorem box}
{\begin{tcolorbox}[colback=yellow!5!white,colframe=yellow!75!black,parbox=false,before upper=\par,before lower=\par]}
{\end{tcolorbox}}
%%%%%%%%%%%%%%%%%%%%%%%%%%%%%%%%%%%%%%%%%%%%%%%%%%%%%%%%%%%%
            
%%%%%%%%%%%%%%%%%%%%% amsthm %%%%%%%%%%%%%%%%
\usepackage{amsthm}

\theoremstyle{definition}
\newtheorem{definition inner}{Definition}[section]
\newenvironment{definition}[1][]
{\begin{definition box}\begin{definition inner}[#1]\hfill\par}
{\end{definition inner}\end{definition box}}

\theoremstyle{plain}
\newtheorem{theorem inner}{Theorem}[section]
\newenvironment{theorem}[1][]
{\begin{theorem box}\begin{theorem inner}[#1]\hfill\par}
{\end{theorem inner}\end{theorem box}}

\theoremstyle{definition}
\newtheorem*{proof inner}{\textit{Proof}}
\renewenvironment{proof}[1][]
{\begin{proof inner}[#1]\hfill\par}
{\qed\end{proof inner}}

\theoremstyle{remark}
\newtheorem*{remark}{\textit{Remark}}
%%%%%%%%%%%%%%%%%%%%%%%%%%%%%%%%%%%%%%%%%%%%%%


\newcommand{\mycases}[2][0.8]{
    \left\{\vphantom{\scalebox{#1}{
        $\begin{matrix}#2\end{matrix}$
    }}\right.\!\!\!\!
    \begin{array}{l@{\quad}l}#2\end{array}
}

\newcommand{\e}{\mathrm{e}{\mkern 1 mu}}

\newcommand{\perm}[2]{\ensuremath{{}_{#1}\mathrm{P}_{#2}}}  % 排列数
\newcommand{\comb}[2]{\ensuremath{{}_{#1}\mathrm{C}_{#2}}}  % 组合数

\newcommand{\E}[1]{\mathbb{E}\qty[#1]}  % 期望
\newcommand{\Prob}{\mathbb{P}\qty}
\newcommand{\Var}{\mathrm{Var}\qty}  % 方差
\newcommand{\Cov}{\mathrm{Cov}\qty}  % 协方差

\begin{document}

\part*{Bivariate Distributions}

\begin{itemize}
    \item a random experiment whose outcome is a scalar
    $\quad\rightarrow\quad$ \textbf{univariate} RV
    \item two random experiments jointly each of whose outcome is a
    scalar,
    
    or a random experiment whose outcome is a pair of
    two scalars $\quad\rightarrow\quad$ \textbf{bivariate} RV
\end{itemize}

\noindent
Let $(X,Y)$ be a bivariate RV. Then
\begin{itemize}
    \item the range: $S\subseteq\mathbb{R}^2$
    \item the range of $X$: $S_X\subseteq\mathbb{R}$,
    and the range of $Y$: $S_Y\subseteq\mathbb{R}$
    \item it holds that $S\subseteq S_X\times S_Y$
\end{itemize}
\begin{remark}
    $X$ and $Y$ can be dependent. Therefore, it may hold that
    $S\subset S_X\times S_Y$.
\end{remark}


\section{Bivariate Distributions of Discrete Type}

\subsection{Definitions}

\begin{definition}[Discrete Bivariate Random Variable]
    A bivariate random variable is said to be of the \textbf{discrete type}
    if its range is \emph{finite} or \emph{countably infinite}.
\end{definition}

\begin{definition}[Joint Probability Mass Function]
    A function $f\colon S\subset\mathbb{R}^2\rightarrow[0,1]$
    is called the \textbf{joint probability mass function} (joint pmf) of
    a DRV $(X,Y)$ with range $S$ if it satisfies:
    \begin{itemize}
        \item $f(x,y)>0$ for $(x,y)\in S$,
        and $f(x,y)=0$ for $(x,y)\notin S$
        \item $\displaystyle\sum_{(x,y)\in S}f(x,y)=1$
        \item $\displaystyle \Prob[(X,Y)\in A]=\sum_{(x,y)\in A}f(x,y)$, where $A\subseteq S$
    \end{itemize}
    Then $S$ is called the \textbf{support} or \textbf{space} of $X$.
\end{definition}

\begin{remark}
    From the third property, we have:
    \begin{equation*}
        \Prob(X=x, Y=y) = f(x, y)
    \end{equation*}
\end{remark}

\begin{remark}
    The summation of the pmf over a set $A\subseteq S$ can be split into two single summations.

    Define $A_X := \{x\mid(x,y)\in A\}$, and
    $A_Y(x) := \{y\mid(x,y)\in A\}$ for $x\in A_X$.
    \begin{equation*}
        \Prob[(X,Y)\in A] =
        \sum_{(x,y)\in A}f(x,y) = \sum_{x\in A_X}\sum_{y\in A_Y(x)}f(x,y)
    \end{equation*}
\end{remark}

\begin{definition}[Marginal Probability Mass Function]
    Let $(X,Y)$ has a joint pmf $f(x,y)$ with space $S$.
    The pmf of $X$ alone,
    called the \textbf{marginal probability mass function}
    (marginal pmf) of $X$, is defined as:
    \begin{equation*}
        f_X(x) := \Prob(X=x) = \sum_{y\in S_Y(x)}f(x,y)
    \end{equation*}
    where
    \begin{equation*}
        S_Y(x) := \{y\mid(x,y)\in S\} \qfor x\in S_X
    \end{equation*}
\end{definition}

\subsection{Independence}

\begin{definition}[Independence]
    $X\in S_X$ and $Y\in S_Y$ are said to be \textbf{independent}
    iff
    \begin{align*}
        \Prob(X=x,Y=y) &= \Prob(X=x)\Prob(Y=y)
    \shortintertext{or equivalently}
        f(x,y) &= f_X(x)f_Y(y)
    \end{align*}
\end{definition}
When $X$ and $Y$ are independent, we have:
\begin{gather*}
    S = S_X \times S_Y \qq{is rectangular} \\
    S_X = S_X(y) \qfor \forall y\in S_Y \\
    S_Y = S_Y(x) \qfor \forall x\in S_X
\end{gather*}
However, the converse is not true.


\subsection{Mathmetical Expectations}

\begin{definition}[Mathmetical Expectation]
    Let $X$ and $Y$ be discrete RVs with joint pmf $f\colon S\rightarrow [0,1]$.
    The \textbf{mathmetical expectation} of a function $g(X,Y)$ is
    \begin{equation*}
        \E{g(X,Y)} := \sum_{(x,y)\in S}g(x,y)f(x,y)
    \end{equation*}
\end{definition}

\begin{theorem}
    Two ways to compute the mean of $X$:
    \begin{equation*}
        \E{X} = \sum_{x\in S_X}x f_X(x)
        = \sum_{(x,y)\in S}x f(x,y)
    \end{equation*}
\end{theorem}


\section{The Corelation Coefficient}

\begin{definition}[Covariance]
    The \textbf{covariance} of two discrete RVs $X$ and $Y$ is defined as
    \begin{equation*}
        \sigma_{X Y} = \Cov(X,Y) := \E{(X-\E{X})(Y-\E{Y})\vphantom{\Big|}}
        = \E{X Y} - \E{X}\E{Y}
    \end{equation*}
    \begin{itemize}
        \item $\Cov(X,Y)<0\quad\implies\quad X$ and $Y$ are \textbf{negatively correlated}
        \item $\Cov(X,Y)>0\quad\implies\quad X$ and $Y$ are \textbf{positively correlated}
        \item $\Cov(X,Y)=0\quad\implies\quad X$ and $Y$ are \textbf{uncorrelated}
    \end{itemize}
\end{definition}

\begin{theorem box}
    \begin{equation*}
        \E{X Y} = \E{X\cdot\E{Y\mid X}}
        = \E{Y\cdot\E{X\mid Y}}
    \end{equation*}
    Meaning: $\E{X Y} = \E{g(X)}$, where $g(x) = x\cdot\E{Y\mid X=x}$.
\end{theorem box}
\begin{proof}
    \begin{equation*}
        \E{X Y} = \sum_{(x,y)\in S} x y f(x,y)
        = \sum_{x\in S_X} x \qty(\sum_{y\in S_y(x)} y\cdot f_{Y|X}(y\mid x))f(x)
        = \sum_{x\in S_X} x\E{Y\mid X=x} f(x)
        = \E{X\cdot\E{Y\mid X=X}}
    \end{equation*}
\end{proof}

\begin{theorem}[Properties of Covariance]
    Let $X,Y,Z$ be RVs, and let $a,b,c,d$ be constants. Then:
    \begin{itemize}
        \item $\Cov(X,X) = \Var(X)$
        \item $\Cov(X,Y) = \Cov(Y,X)$
        \item $\Cov(X+Y,Z) = \Cov(X,Z) + \Cov(Y,Z)$
        \item $\Cov(a X + b, c Y + d) = a c\ \Cov(X,Y)$
        \item $\Cov(X, a) = 0$
    \end{itemize}
\end{theorem}

\begin{theorem}
    In terms of two RVs $X$ and $Y$, generally,
    \begin{align*}
        \qq{Independent} &\implies \qq{Uncorrelated}\qq{i.e.}\Cov(X,Y)=0 \\
        \qq{Independent} &\centernot\impliedby \qq{Uncorrelated}
    \end{align*}
\end{theorem}
\begin{proof}[Independence implies uncorrelation]
    \begin{equation*}
        f(x,y) = f_X(x)f_Y(y) \quad\implies\quad
        \E{X Y} = \sum_{x\in S_X}\sum_{y\in S_Y} x y f_X(x)f_Y(y)
        = \E{X}\E{Y} \quad\implies\quad
        \Cov(X,Y) = 0
    \end{equation*}
\end{proof}

\begin{definition}[Correlation Coefficient]
    The \textbf{correlation coefficient} of two discrete RVs $X$ and $Y$ is defined as
    \begin{equation*}
        \rho(X,Y) := \frac{\Cov(X,Y)}{\sqrt{\Var(X)\Var(Y)}}
        = \frac{\sigma_{X Y}}{\sigma_X\sigma_Y}
        \in [-1,1]
    \end{equation*}
    provided that $\sigma_X\sigma_Y\neq 0$.
\end{definition}

\begin{remark}
    $\rho$ is a normalized  version of $\Cov(X,Y)$.
\end{remark}

\begin{theorem}
    $\rho$ provides a normalized measure of the strength of the
    linear relationship between $X$ and $Y$.
    \begin{equation*}
        \abs{\rho} = 1 \quad\iff\quad
        \exists c \qq{s.t.}
            Y - \E{Y} = c(X-\E{X})
    \end{equation*}
    where $c$ shares the same sign as $\rho$ and
    \begin{equation*}
        \abs{c} = \sqrt{\frac{\Var(Y)}{\Var(X)}}
    \end{equation*}
\end{theorem}

\begin{theorem}[Cauchy-Schwarz Inequality]
    For any two RVs $X$ and $Y$,
    \begin{equation*}
        \Cov(X,Y)^2 \leqslant \Var(X)\Var(Y)
    \end{equation*}
\end{theorem}
\begin{proof}[Cauchy-Schwarz Inequality]
    Let $U=X-\E{X}$ and $V=Y-\E{Y}$.
    Let $f(t)=\E{(U+t V)^2}$. Then
    \begin{equation*}
        f(t) = \E{U^2} + 2 t\E{U V} + t^2\E{V^2}
        = \Var(X) + 2 t\Cov(X,Y) + t^2\Var(Y)
        \geqslant 0
    \end{equation*}
    The discriminant of $f(t)$ is
    \begin{equation*}
        \Delta = 4\Cov(X,Y)^2 - 4\Var(X)\Var(Y)
        = 4\Cov(X,Y)^2 - 4\Var(X)\Var(Y) \leqslant 0
    \end{equation*}
\end{proof}

\section{Conditional Distributions}

\begin{definition}[Conditional Probability Mass Function]
    The \textbf{conditional pmf} of $X$ given $Y=y$ is defined as
    \begin{equation*}
        f_{X|Y}(x\mid y) := \frac{f_{X,Y}(x,y)}{f_Y(y)}
        \qfor x\in S_X(y)
    \end{equation*}
    provided that $f_Y(y)>0$.
\end{definition}

\begin{proof}
    \begin{equation*}
        f_{X|Y}(x\mid y) = \Prob(X=x\mid Y=y) = \frac{\Prob(X=x,Y=y)}{\Prob(Y=y)}
        = \frac{f_{X,Y}(x,y)}{f_Y(y)}
    \end{equation*}
\end{proof}

\begin{remark}
    Easy to check that conditional pmf is well-defined.
\end{remark}
\begin{remark}
    When $X$ and $Y$ are independent, we have:
    \begin{equation*}
        f_{X|Y}(x\mid y) = \frac{f_X(x)f_Y(y)}{f_Y(y)} = f_X(x)
    \end{equation*}
\end{remark}
    
\begin{definition}[Conditional Mathmetical Expectation]
    Let $g(Y)$ be a function of $Y$. The \textbf{conditional expectation} of
    $g(Y)$ given $X=x$ is defined as
    \begin{equation*}
        \E{g(Y)\mid X=x} := \sum_{y\in S_Y(x)}g(y)\cdot f_{Y|X}(y\mid x)
    \end{equation*}
    where $x\in S_X$.
    \begin{itemize}
        \item Conditional mean:
        $\quad \E{Y\mid X=x}$
        \item Conditional variance:
        $\quad \Var(Y\mid X=x) := \E{(Y-\E{Y\mid X=x})^2\mid X=x}$
    \end{itemize}
\end{definition}
\begin{remark}
    \begin{equation*}
        \Var(Y\mid X=x) = \E{Y^2\mid X=x} - \big(\E{Y\mid X=x}\big)^2
    \end{equation*}
\end{remark}

\begin{theorem}
    \begin{equation*}
        \E{X} = \E{\E{X\mid Y=Y}} = \sum_{y\in S_Y}\E{X\mid Y=y}f_Y(y)
        = \sum_{(x,y)\in S} x f_{X,Y}(x,y)
    \end{equation*}
\end{theorem}


\section{Examples}

\subsection{Trinomial Distribution}

\paragraph{Desciption}
A random experiment that has three \emph{mutually exclusive}
and \emph{exhaustive} outcomes is repeated $n$ independent times,
and the probabilities of each outcome remain the same
for each repetition.
Such $n$ repetitions can be called a trinomial experiment.
Let $X$, $Y$, and $n-X-Y$ be the number of the three outcomes,
respectively.

\begin{definition}[Trinomial Distribution]
    \textbf{Parameters}\qquad $n=1,2,3,\cdots \qquad p_X\in[0,1]
    \qquad p_Y\in[0,1] \qq{s.t.} p_X+p_Y\leqslant 1$.

    $(X,Y)$ is said to have a \textbf{trinomial distribution},
    if the \textbf{joint pmf} is given by
    \begin{equation*}
        f_{X,Y}(x,y) = \binom{n}{x,y,n-x-y}p_X^x p_Y^y (1-p_X-p_Y)^{n-x-y}
        = \frac{n!}{x!y!(n-x-y)!}p_X^x p_Y^y (1-p_X-p_Y)^{n-x-y}
    \end{equation*}
    The \textbf{marginal pmf} of $X$ is given by
    \begin{equation*}
        f_X(x) = \binom{n}{x}p_X^x(1-p_X)^{n-x}
    \end{equation*}
    Namely, $X\sim\mathrm{B}(n,p_X)$, $Y\sim\mathrm{B}(n,p_Y)$.
\end{definition}

\begin{remark}
    Trinomial Expansion
    \begin{equation*}
        (a+b+c)^n = \sum_{x=0}^{n}\sum_{y=0}^{n-x}
        \binom{n}{x,y,n-x-y}a^x b^y c^{n-x-y}
        = \sum_{x=0}^{n}\sum_{y=0}^{n-x}
        \frac{n!}{x!y!(n-x-y)!}a^x b^y c^{n-x-y}
    \end{equation*}
\end{remark}

\begin{theorem}[Trinomial Distribution]
    The conditional pmf of $Y$ given $X=x$ is
    \begin{equation*}
        f_{Y|X}(y\mid x) =
        \binom{n-x}{y} \qty(\frac{p_Y}{1-p_X})^y \qty(1-\frac{p_Y}{1-p_X})^{n-x-y}
    \end{equation*}
    Namely, $(Y\mid X=x)\sim\mathrm{B}\qty(n-x,\dfrac{p_Y}{1-p_X})$,
    and $(X\mid Y=y)\sim\mathrm{B}\qty(n-y,\dfrac{p_X}{1-p_Y})$.

    The covariance of $X$ and $Y$ can be computed as
    \begin{equation*}
        \E{X Y} = n(n-1)\cdot p_X p_Y \qquad\qquad
        \Cov(X,Y) = -n \cdot p_X p_Y
    \end{equation*}
\end{theorem}
\begin{proof}[Covariance of trinomial distribution]
    \noindent\textbf{Method 1} \quad Conditional Expectation Method
    \begin{gather*}
        \E{Y^2} = (n^2-n)\cdot p_Y^2 + n\cdot p_Y \\
        \E{X Y} = \E{Y\cdot\E{X\mid Y}}
        = \frac{p_X}{1-p_Y}\E{Y(n-Y)}
        = n(n-1)\cdot p_X p_Y
    \end{gather*}
    \noindent\textbf{Method 2} \quad Indicator Variable Method
    \par
    For each trial $i$, define the indicator variables $X_i$ and $Y_i$:
    \begin{equation*}
        X_i = \mycases{
            1 & \text{if the $i$-th trial is the outcome of $X$} \\
            0 & \text{otherwise}
        } \qquad\qquad
        Y_i = \mycases{
            1 & \text{if the $i$-th trial is the outcome of $Y$} \\
            0 & \text{otherwise}
        }
    \end{equation*}
    Then $X_i$ and $Y_j$ are independent for $i\neq j$,
    and $X_i Y_i=0$. Therefore,
    \begin{align*}
        \E{X Y} = \E{\sum_{i=1}^n X_i \sum_{i=1}^n Y_i}
        = \sum_{i=1}^n \sum_{j=1}^n \E{X_i Y_j}
        = \sum_{i\neq j}\E{X_i}\E{Y_j}
        = \sum_{i\neq j} p_X p_Y
        = n(n-1) p_X p_Y
    \end{align*}
\end{proof}


\section{Bivariate Distributions of Continuous Type}

\begin{definition}[Joint Probability Density Function]
    A function $f\colon S\subseteq\mathbb{R}^2\rightarrow[0,+\infty)$
    is called the \textbf{joint probability density function} (joint pdf) of
    a CRV $(X,Y)$ with range $S$ if it satisfies:
    \begin{itemize}
        \item $f(x,y)>0$ for $(x,y)\in S$,
        and $f(x,y)=0$ for $(x,y)\notin S$
        \item $\displaystyle\iint_{S}f(x,y)\dd{x}\dd{y}=1$
        \item $\displaystyle \Prob[(X,Y)\in A]=\iint_{A}f(x,y)\dd{x}\dd{y}$,
        where $A\subseteq S$
    \end{itemize}
    Then $S$ is called the \textbf{support} or \textbf{space} of $X$.
\end{definition}

\begin{definition}[Marginal Probability density Function]
    Let $(X,Y)$ has a joint pdf $f(x,y)$ with space $S$.
    The pdf of $X$ alone,
    called the \textbf{marginal probability density function}
    (marginal pdf) of $X$, is defined as:
    \begin{equation*}
        f_X(x) :=  \int_{S_Y(x)}\!\!\!\!\!\!\!\!\!f(x,y)\dd{y}
        \qfor x\in S_X
    \end{equation*}
    where
    \begin{equation*}
        S_Y(x) := \{y\mid(x,y)\in S\} \qfor x\in S_X
    \end{equation*}
\end{definition}


\subsection{Mathmetical Expectations}

\begin{definition}[Mathmetical Expectation]
    Let $X$ and $Y$ be CRVs with joint pdf $f\colon S\rightarrow [0,+\infty)$.
    The \textbf{mathmetical expectation} of a function $g(X,Y)$ is
    \begin{equation*}
        \E{g(X,Y)} := \iint_{S}g(x,y)f(x,y)\dd{x}\dd{y}
    \end{equation*}
\end{definition}
\begin{theorem}
    \begin{equation*}
        \E{X} = \int_{S_X}x f_X(x)\dd{x}
        = \iint_{S}x f(x,y)\dd{x}\dd{y}
        = \int_{S_Y}\!\!\!\! \E{X\mid Y=y}f_Y(y)\dd{y}
        = \E{\E{X\mid Y=Y}}
    \end{equation*}
\end{theorem}

\subsection{Covariance \& Correlation Coefficient}

\begin{definition}[Independence]
    $X\in S_X$ and $Y\in S_Y$ are said to be \textbf{independent}
    iff
    \begin{align*}
        f(x,y) &= f_X(x)f_Y(y)
    \end{align*}
\end{definition}

If $X$ and $Y$ are independent, then
\begin{gather*}
    S = S_X\times S_Y
\end{gather*}

\begin{definition}[Covariance \& Correlation Coefficient]
    Let $X$ and $Y$ be CRVs.
    \begin{align*}
        \Cov(X,Y) &:= \E{(X-\E{X})(Y-\E{Y})\vphantom{\Big|}} \\
        &= \E{X Y} - \E{X}\E{Y}
    \end{align*}
    \begin{gather*}
        \rho(X,Y) := \frac{\Cov(X,Y)}{\sqrt{\Var(X)\Var(Y)}}
        \qfor \Var(X)\Var(Y) \neq 0
    \end{gather*}
\end{definition}

\subsection{Conditional Distributions}


\begin{definition}[Conditional Probability Density Function]
    The \textbf{conditional pdf} of $X$ given $Y=y$ is defined as
    \begin{equation*}
        f_{X|Y}(x\mid y) := \frac{f_{X,Y}(x,y)}{f_Y(y)}
        \qfor x\in S_X(y)
    \end{equation*}
    provided that $f_Y(y)>0$.
\end{definition}

\begin{definition}[Conditional Mathmetical Expectation]
    Let $g(Y)$ be a function of $Y$. The \textbf{conditional expectation} of
    $g(Y)$ given $X=x$ is defined as
    \begin{equation*}
        \E{g(Y)\mid X=x} := \int_{S_Y(x)}\!\!\!\!\!\!\!\!\!g(y)f_{Y|X}(y\mid x)\dd{y}
    \end{equation*}
    where $x\in S_X$.
    \begin{itemize}
        \item Conditional mean:
            $\quad\displaystyle
            \E{Y\mid X=x} = \int_{S_Y(x)}\!\!\!\!\!\!\!\!\!y\cdot f_{Y|X}(y\mid x)\dd{y}$
            \item Conditional variance:
            $\quad \begin{aligned}[t]
                \Var(Y\mid X=x) &:= \E{(Y-\E{Y\mid X=x})^2\mid X=x} \\
                &= \E{Y^2\mid X=x} - \big(\E{Y\mid X=x}\big)^2
            \end{aligned}$
    \end{itemize}
\end{definition}


\section{Bivariate Normal Distribution}

\begin{definition}[Bivariate Normal Distribution]
    \textbf{Parameters}\qquad $\mu_X\in\mathbb{R} \qquad \mu_Y\in\mathbb{R} \qquad
    \sigma_X>0 \qquad \sigma_Y>0 \qquad \rho\in[-1,1]$

    $X$ and $Y$ are said to have a \textbf{bivariate normal distribution} if the joint pdf is given by
    \begin{gather*}
        f(x,y) = \frac{1}{2\pi\sigma_X\sigma_Y\sqrt{1-\rho^2}}
        \exp(-\frac{1}{2(1-\rho^2)}\qty[
            \qty(\frac{x-\mu_X}{\sigma_X})^2
            - 2\rho\qty(\frac{x-\mu_X}{\sigma_X})\qty(\frac{y-\mu_Y}{\sigma_Y})
            +\qty(\frac{y-\mu_Y}{\sigma_Y})^2
        ])
    \end{gather*}
    where $\rho$ is the \emph{correlation coefficient} of $X$ and $Y$.
\end{definition}

\begin{theorem}[Bivariate Normal Distribution]
    \begin{itemize}
        \item Marginal:\quad
            $\displaystyle
            X\sim\mathrm{N}(\mu_X,\sigma_X^2),\quad
            Y\sim\mathrm{N}(\mu_Y,\sigma_Y^2)$
        \item Conditional:\quad
        $\begin{gathered}[t]
            [X\mid Y=y] \sim\mathrm{N}\qty(
                \mu_X + \rho\frac{\sigma_X}{\sigma_Y}(y-\mu_Y),\;
                (1-\rho^2)\sigma_X^2
            ) \\
            [Y\mid X=x] \sim\mathrm{N}\qty(
                \mu_Y + \rho\frac{\sigma_Y}{\sigma_X}(x-\mu_X),\;
                (1-\rho^2)\sigma_Y^2
            )
        \end{gathered}$
        \item Covariance:\quad
            $\Cov(X,Y) = \rho\sigma_X\sigma_Y$
        \item Correlation Coefficient:\quad
            $\rho(X,Y) = \rho$
    \end{itemize}
\end{theorem}

\begin{proof}[Covariance of bivariate normal distribution]
    Since $X\sim\mathrm{N}(\mu_X,\sigma_X^2)$, $\E{X^2}=\sigma_X^2-\mu_X^2$.
    Moreover, $\E{Y\mid X=X} = \mu_Y + \rho\dfrac{\sigma_Y}{\sigma_X}(x-\mu_X)$.
    Then,
    \begin{gather*}
        \E{X Y} = \E{X\cdot\E{Y\mid X=X}\vphantom{\Big|}}
        = \rho\frac{\sigma_X}{\sigma_Y}\E{X^2}
        + \qty(\mu_Y-\mu_X\rho\frac{\sigma_X}{\sigma_Y})\E{X}
        = \rho\sigma_X\sigma_Y + \mu_X\mu_Y \\
        \Cov(X,Y) = \E{X Y} - \E{X}\E{Y} = \rho\sigma_X\sigma_Y
    \end{gather*}
\end{proof}

\begin{theorem}
    For a bivariate normal distribution,
    \begin{equation*}
        \qq{Independence} \iff \qq{Uncorrelation}
    \end{equation*}
\end{theorem}
\begin{proof}
    For any bivariate RV,
    independence implies uncorrelation.
    Therefore, we only need to prove the converse.
    \begin{gather*}
        \Cov(X,Y) = 0 \quad\implies\quad \rho = 0 \quad\implies\quad
        [X\mid Y=y] \sim\mathrm{N}(\mu_X,\sigma_X^2)
        \quad\implies\quad [X\mid Y=y] = X
    \end{gather*}
\end{proof}

\end{document}
